\subsection{Matched Precision, Work, and Anchor Demand}
\label{sec:r7_matched_evidence}
We now hold fixed what differed in the retained global comparison. For every one of the 51 original anchor intervals, the corrected chord, localized count recursion, and rectangular upper recursion use the same two endpoint policies and the same eight closed Bernstein cells. Every anchor is freshly solved on the full $1,617$-state, eight-date, $1,568$-action target, and its policy coefficients are independently reevaluated. The shared values and coefficients reproduce their R6 records to the deposited numerical tolerance. The three full-bank rows therefore reproduce the sixth-round matched countercomparison while adding a common-machine work account. Two additional $d=0.5$ banks retain every second or fourth anchor and include the right endpoint.
\begin{table}[!htbp]
\centering\small
\caption{Matched Upper-Oracle Precision and Work}
\label{tab:r7_matched}
\begin{tabular}{rrlrrr}
\toprule
$d$ & Anchors & Upper oracle & Bound & Upper sec. & Total sec. \\
\midrule
0 & 19 & Compressed & 0.000904573 & 27.02 & 53.34 \\
 &  & Local count & 0.000733846 & 95.12 & 121.40 \\
 &  & Rectangular & 0.0187608 & 21.62 & 47.90 \\
0.5 & 18 & Compressed & 0.000899829 & 25.46 & 51.01 \\
 &  & Local count & 0.000685103 & 88.91 & 114.43 \\
 &  & Rectangular & 0.0171499 & 20.01 & 45.52 \\
0.5 & 6 & Compressed & 0.00698458 & 7.34 & 19.26 \\
 &  & Local count & 0.00338132 & 25.61 & 37.52 \\
 &  & Rectangular & 0.0427253 & 5.75 & 17.66 \\
0.5 & 10 & Compressed & 0.00327336 & 13.35 & 29.79 \\
 &  & Local count & 0.0020246 & 46.32 & 62.74 \\
 &  & Rectangular & 0.0255638 & 10.35 & 26.77 \\
1 & 17 & Compressed & 0.000941412 & 23.64 & 47.04 \\
 &  & Local count & 0.000435745 & 81.90 & 105.27 \\
 &  & Rectangular & 0.0150829 & 18.24 & 41.60 \\
\bottomrule
\end{tabular}
\par\smallskip
\begin{minipage}{0.98\linewidth}\footnotesize
Notes: All states, dates, and $\lambda\in[0,1]$; identical two-policy lower banks and eight cells within each compared interval. The 19-, 18-, and 17-anchor banks are the original R6 banks; the 6- and 10-anchor banks coarsen the $d=0.5$ bank. Total time includes common kernel setup, anchor solves, policy evaluation, coefficient restriction, upper construction, and certification. Shared stages are attributed equally. One serial single-thread pass; diagnostic replay is excluded equally. Kernel-call and persistent-array counts, every interval bound, and stage times are in the replication records.
\end{minipage}
\end{table}


The localized count certificate is sharper in every full-bank row, as Proposition \ref{prop:r7_dominance} requires. The compressed construction buys smaller upper-oracle work rather than smaller loss. At eight dates its 26 action-wide kernel applications per interval compare with 70 for the cached count recursion and 16 for the rectangular relaxation. The persistent correction needs seven additional state vectors, compared with 28 interior count vectors; neither number includes the shared endpoint values or the uncompressed lower-policy bank. The rectangular relaxation is cheaper in kernel work but loses much more precision here by allowing repeated resets of the regime probability. Thus the table separates a validity result, a sharpness result, and a cost result.

All elapsed-time columns include kernel construction, the required anchor solves, policy evaluation, local coefficient restriction, upper construction, and subdivision. Shared stages are measured once and attributed equally. The compact implementation expands one date's upper coefficients during certification; this work is included rather than removed from the timing. These are one-pass serial single-thread measurements, not medians of independently rebuilt machines. The operation counts provide a reproducible work comparison independent of elapsed-time noise. No new neural training or learned-speed advantage is inferred from them.

Common banks do not determine how many anchors an independently refined method would request. Table \ref{tab:r7_adaptive} therefore starts three procedures at $[0,1]$ with empty per-arm caches, fixes $d=0.5$, and uses the same left-first dyadic subdivision and $10^{-3}$ target. The chord-only and count-only arms use their own certificates to decide whether to split. The cascade tries count only after the chord fails. Every rejected parent interval, newly solved anchor, lower evaluation, and cell test is included in that arm's cost. This is an autonomous refinement experiment, unlike the coarsened-bank rows of Table \ref{tab:r7_matched}.
\begin{table}[!htbp]
\centering\small
\caption{Autonomous Refinement from Empty Policy Caches}
\label{tab:r7_adaptive}
\begin{tabular}{lrrrr}
\toprule
Rule & Anchors & Final bound & Kernel calls & Total sec. \\
\midrule
Chord & 18 & 0.000899829 & 858 & 73.41 \\
Count & 11 & 0.000924471 & 1330 & 115.38 \\
Cascade & 11 & 0.000924471 & 1614 & 127.73 \\
\bottomrule
\end{tabular}
\par\smallskip
\begin{minipage}{0.98\linewidth}\footnotesize
Notes: Full target, $d=0.5$, tolerance $10^{-3}$. Each arm starts at $[0,1]$ and includes all rejected parent intervals. Kernel calls are action-wide upper-oracle applications, not anchor-solve calls. Total time includes both. Cascade tries local count only after compressed certification fails. All arms use identical dyadic split order, two endpoint policies, and eight certificate cells.
\end{minipage}
\end{table}

The chord-only procedure requests 18 anchors; local count requests 11; the cascade requests 11. Thus the cascade avoids 7 endpoint optimizations relative to chord-only refinement at the same certified tolerance. Its elapsed time is 127.73 seconds, against 73.41 and 115.38 for the two single-oracle procedures, respectively. The measured costs, rather than the anchor reduction alone, determine the computational comparison.

The cascade's anchor set is checked to be a subset of the chord-only set, as the nesting argument predicts. It need not be the fastest procedure: buying a sharper certificate also consumes work. These finite-model results support a choice among upper-information budgets; they do not establish a dimension-free neural advantage or the optimality of any particular refinement strategy.

\subsection{A Certified Preference-Adjustment Effect on a Region}
\label{sec:r7_region_evidence}
The economic certificate holds the operating boundary, settlement function, consumption menu, financial technology, and stochastic preference shocks fixed. At $(u,X)=(2,1.25)$ and $k=2$, consider the whole rectangle
\[
 (\lambda,d)\in[0,0.25]\times[0.4,0.45].
\]
The parameter $\lambda$ is the same per-period probability of the positive-co-movement quadrature regime used above, not a new reward coefficient. Each risk class is solved at the four corners with and without deliberate adjustment. Their feasible policies are then evaluated as transition polynomials with affine contract coefficients. Localized count upper values at the contract endpoints are interpolated in $d$ using convexity of the optimal class value. Proposition \ref{prop:r7_tensor} certifies the resulting tensor coefficients on the entire rectangle.
\begin{table}[!htbp]
\centering\small
\caption{A Joint-Region Certificate for the Adjustment Effect}
\label{tab:r7_decision}
\begin{tabular}{lrrr}
\toprule
Adjustment regime & Lower $\Delta$ & Upper $\Delta$ & Max. class gap \\
\midrule
Adjustment available & $6.88873\times10^{-5}$ & 0.000324537 & $7.5953\times10^{-9}$ \\
No deliberate adjustment & -0.000369485 & $-5.29937\times10^{-5}$ & $2.19355\times10^{-6}$ \\
\bottomrule
\end{tabular}
\par\smallskip
\begin{minipage}{0.98\linewidth}\footnotesize
Notes: Initial state $(u,X)=(2,1.25)$, $k=2$, full finite target, entire $(\lambda,d)\in[0,0.25]\times[0.4,0.45]$. Bounds include the arithmetic allowance. Maximum class gap is the larger of the two class-specific upper-minus-feasible-lower certificates in that regime. The per-class arithmetic allowance is $3.79765\times10^{-9}$; the lower bound on the relative adjustment option is 0.000121881. Eight corner optimizations produce sixteen feasible class policies; eight count constructions certify the rectangle. No generic $10^{-3}$ welfare tolerance is used for these signs.
\end{minipage}
\end{table}


The positive first-risk class is strictly preferred throughout the region when adjustment is available, and the nonpositive class is strictly preferred throughout the same region when it is unavailable. Both signs survive the separate numerical allowance. Their difference bounds the relative adjustment option from below, so the conclusion identifies an incremental preference-formation effect under unchanged financial and settlement primitives. It is not the assertion that adjustment is necessary for every duration-induced risk reversal. The earlier $d=1$ comparison, where both regimes select positive risk, remains in the paper.

This calculation uses the transition representation to certify the economically relevant decision. The upper-information and feasible-policy errors are measured separately for the four class values; they are not assigned the generic $10^{-3}$ welfare tolerance. The tensor coefficients and a coefficient-only replay accompany the policy bank. Supplement S.11 gives the round-to-nearest arithmetic account relative to the stored finite model. Independent corner replay is an implementation check, not a replacement for the region inequality. The result concerns this initial state and rectangle, not all initial states, uniqueness of a global switching threshold, or a continuous-time discretization theorem.
